\chapter{Cast Application}

\section*{Overview}
A cast circumferentially immobilizes an injured limb (usually a fracture) using plaster or fiberglass bandages over cotton padding. This maintains bone alignment during healing.

\section*{Alternative Names}
Orthopedic casting, Plaster cast

\section*{Principle}
Wet plaster or fiberglass is wrapped around the limb and molded to its contours, then hardens into a rigid shell. The cast immobilizes the bone ends or joint, holds the fracture in correct alignment, and prevents the movement that would disrupt healing.
\section*{History}
Antonius Mathijsen, a Dutch military surgeon, invented plaster-of-Paris bandages in 1851. Fiberglass casts were introduced in the 1970s.

\section*{Why It Is Done}
Ordered to immobilize fractures, severe sprains, or post-surgical repairs to prevent movement and facilitate bone and tissue healing.

\section*{Is This a Primary or Secondary Test or Procedure}
Primary conservative treatment for stable, reduced bone fractures.

\section*{How It Is Performed}
A stockinette is applied to the limb, followed by layers of soft cotton padding. Plaster or fiberglass rolls are wet and wrapped around the limb, then molded by the clinician.

\section*{Preparation}
Clean and dry the skin. Ensure the fracture is adequately reduced and aligned.

\section*{Contraindications and Precautions}
Severe acute swelling (relative; a splint should be used initially to avoid compartment syndrome).

\section*{Risks}
Compartment syndrome (pressure buildup), skin pressure ulcers/necrosis, joint stiffness, or thermal burns during plaster setting.

\section*{Aftercare and Recovery}
Elevate the limb above the heart for 24-48 hours. Monitor for '5 Ps' of compartment syndrome: Pain, Pulselessness, Pallor, Paresthesia, Paralysis. Keep the cast dry.

\section*{Outcome and Follow-up}
Signs of a successful application include bone stability, comfortable fit, and distal limb warmth and pink color.

\section*{Expected Results}
Rigid immobilization of the joint/bone; warm distal extremities; capillary refill <2 seconds; intact sensation.

\section*{Follow-up}
Cast check in 1 week; serial X-rays to ensure stable alignment; cast removal in 4-8 weeks.

\section*{References}
\begin{enumerate}
  \item MedlinePlus. Cast Application. U.S. National Library of Medicine; 2025.
  \item Current Medical Diagnosis and Treatment. McGraw-Hill; 2025.
\end{enumerate}

\clearpage
